\documentclass{article}

\usepackage{amsmath,amssymb}
\usepackage{xurl}
\usepackage[hidelinks]{hyperref}
\setlength{\emergencystretch}{2em}

% Shared commands for notes in docs/paper-gaps/.

\DeclareUnicodeCharacter{2080}{\ifmmode{}_{0}\else\textsubscript{0}\fi}
\DeclareUnicodeCharacter{2081}{\ifmmode{}_{1}\else\textsubscript{1}\fi}
\DeclareUnicodeCharacter{2082}{\ifmmode{}_{2}\else\textsubscript{2}\fi}
\DeclareUnicodeCharacter{2099}{\ifmmode{}_{n}\else\textsubscript{n}\fi}

\newcommand{\Lean}{\textsc{Lean}}
\ExplSyntaxOn
\NewDocumentCommand{\leanid}{m}{
  \ifmmode
    \textsf{\detokenize{#1}}
  \else
    \group_begin:
    \urlstyle{sf}
    \tl_set:Nn \l_tmpa_tl {#1}
    \tl_replace_all:Nnn \l_tmpa_tl {\_} {_}
    \exp_args:NV \path \l_tmpa_tl
    \group_end:
  \fi
}
\ExplSyntaxOff
\providecommand{\tr}{\operatorname{tr}}
\newcommand{\tnleanIssueBase}{https://github.com/LionSR/TNLean/issues/}
\newcommand{\tnleanPrBase}{https://github.com/LionSR/TNLean/pull/}
\newcommand{\tnleanDocBase}{https://sirui-lu.com/TNLean/docs/}
\newcommand{\ghissue}[1]{\href{\tnleanIssueBase#1}{GitHub issue \##1}}
\newcommand{\ghpr}[1]{\href{\tnleanPrBase#1}{PR \##1}}
\newcommand{\doclink}[1]{\href{\tnleanDocBase#1.html}{\path{#1}}}

% Dirac notation and the identity operator, as in blueprint/src/macros/common.tex.
% \providecommand keeps note-local definitions authoritative.
\usepackage{dsfont}
\providecommand{\ket}[1]{|#1\rangle}
\providecommand{\bra}[1]{\langle #1|}
\providecommand{\braket}[2]{\langle #1 | #2 \rangle}
\providecommand{\ketbra}[2]{|#1\rangle\!\langle #2|}
\providecommand{\Id}{\mathds{1}}
\providecommand{\one}{\mathds{1}}
\providecommand{\C}{\mathbb{C}}
\providecommand{\MN}[1]{M_{#1}(\C)}
\providecommand{\MD}{M_D(\C)}

% External referencing for paper sources.  The build helper writes sanitized
% auxiliary files under build/paper-aux/ and excludes duplicated source labels.
\usepackage{xr}

% Import a generated paper auxiliary file with a caller-supplied label prefix.
% The candidate paths support builds from docs/paper-gaps/, the repository root,
% and build/paper-aux/ itself.
\newcommand{\importpaperaux}[2]{%
  \IfFileExists{../../build/paper-aux/#2.aux}{%
    \externaldocument[#1]{../../build/paper-aux/#2}%
  }{%
    \IfFileExists{build/paper-aux/#2.aux}{%
      \externaldocument[#1]{build/paper-aux/#2}%
    }{%
      \IfFileExists{#2.aux}{%
        \externaldocument[#1]{#2}%
      }{}%
    }%
  }%
}

% General external reference macros
\newcommand{\paperref}[2]{\ref{#1-#2}}
\newcommand{\papereqref}[2]{(\ref{#1-#2})}

% CPSV16 source (1606.00608 / MPDO-22-12-17-2.tex) external document and shortcuts
\importpaperaux{cpsv16-}{1606.00608/MPDO-22-12-17-2-tnlean}
\newcommand{\cpsvref}[1]{\paperref{cpsv16}{#1}}
\newcommand{\cpsveqref}[1]{\papereqref{cpsv16}{#1}}

% Verdict marker: \gapnote{<kind>}{<status>}. Kind is the policy
% classification (clarification, local-correction, scope-restriction,
% unfaithful, false-source, open-gap); status is open, wip (elimination
% underway), resolved, or historical. Machine-read by the site index;
% prints a verdict line.
\newcommand{\gapnote}[2]{\par\noindent\textbf{Verdict:} #1 (#2).\par}


\title{Wolf Chapter~5 Operator Jensen and Lieb Concavity Axioms}
\author{}
\date{2026-06-22}

\begin{document}
\maketitle

\section{Scope}

This note records the mathematical content of the four axioms that remain in the
Chapter~5 operator-convexity development\footnote{%
    The axioms are
    \leanid{posMap_rpow_concave_jensen}, \leanid{posMap_rpow_convex_jensen},
    \leanid{posMap_log_concave_jensen}, and \leanid{lieb_concavity_axiom} in
    \doclink{TNLean/Axioms/OperatorConvexity}.}
together with the cited sources they stand in for, the precise gap that blocks a
proof, and the first concrete step toward eliminating each gap.  The development
follows \cite[Theorem~5.1]{Wolf2012Quantum}.

The notation is fixed throughout.  Operators act on a finite-dimensional Hilbert
space, identified with $D\times D$ complex matrices.  The order $A\le B$ is the
Loewner order ($B-A$ positive semidefinite).  A linear map $T$ on matrices is
\emph{positive} if it sends positive semidefinite operators to positive
semidefinite operators, \emph{subunital} if $T(\mathbf 1)\le\mathbf 1$, and
\emph{unital} if $T(\mathbf 1)=\mathbf 1$.  The power $A^p$ for $A\ge 0$ and the
logarithm $\log A$ for $A>0$ are defined by the continuous functional calculus.

\section{The three positive-map Jensen axioms}

\subsection{The assertions}

For a positive subunital map $T$ the development asserts the operator Jensen
inequalities
\begin{align}
  T(A^p) &\le (T A)^p, & p&\in[0,1], \ A\ge 0, \tag{\path{posMap_rpow_concave_jensen}}\\
  (T A)^p &\le T(A^p), & p&\in[1,2], \ A\ge 0, \tag{\path{posMap_rpow_convex_jensen}}
\end{align}
and, for a positive \emph{unital} map $T$ and $A>0$,
\begin{align}
  T(\log A) &\le \log(T A). \tag{\path{posMap_log_concave_jensen}}
\end{align}
These are the operator-valued Jensen inequalities behind Wolf
Corollary~5.2.\footnote{%
    They are consumed by
    \leanid{cor52_item1_rpow_of_subunital},
    \leanid{cor52_item2_rpow_of_subunital}, and
    \leanid{cor52_item3_log_of_subunital} in
    \doclink{TNLean/Channel/Schwarz/OperatorMonotone}.}
The first and third are concavity inequalities (for the operator-concave
functions $x\mapsto x^p$ on $[0,1]$ and $x\mapsto\log x$); the second is a
convexity inequality (for the operator-convex function $x\mapsto x^p$ on
$[1,2]$).  The logarithmic case requires $T$ unital, not merely subunital: with
$T$ only subunital the inequality fails because $\log$ is unbounded below near
$0$.

\subsection{The point at issue: the genuine gap is positive-map Jensen}

Earlier notes on this issue recorded the obstruction as the absence of the
Loewner integral representation of $x^p$ from the formal library.  That
diagnosis is now superseded.  Mathlib~4.31 supplies the integral
representation\footnote{%
    \leanid{CFC.exists_measure_nnrpow_eq_integral_cfcₙ_rpowIntegrand₀₁}
    and its $[1,2]$ companion in
    \path{Mathlib/Analysis/SpecialFunctions/ContinuousFunctionalCalculus/Rpow/IntegralRepresentation.lean},
    together with the Bochner-integral convexity bridge for C\textsuperscript{*}-algebras.}
and, through it, the operator concavity of $x\mapsto x^p$ on $[0,1]$
(\leanid{CFC.concaveOn_rpow}) and of $\log$ (\leanid{CFC.concaveOn_log}).  The
scalar inputs of the two concavity axioms are therefore already
operator-concave in the formal library.  The scalar input of the convex axiom,
operator convexity of $x\mapsto x^p$ on $[1,2]$, is \emph{not} yet in the
library: only the operator-monotone and operator-concave $[0,1]$ lemmas exist,
and operator convexity of \texttt{rpow} over $\mathrm{Icc}\,1\,2$ is an open
upstream TODO.\footnote{%
    \path{Mathlib/Analysis/SpecialFunctions/ContinuousFunctionalCalculus/Rpow/Order.lean},
    line~29.}
That missing convexity lemma is, however, a small derivable increment
(\S\ref{ssec:first-step}), not the deep obstruction.

The remaining gap is the step \emph{from} operator concavity (a statement about a
single operator) \emph{to} the Jensen inequality $T(fA)\le f(TA)$ for a positive
map.  This is the Choi--Davis--Jensen inequality (the positive-map form of the
Hansen--Pedersen operator Jensen inequality): for an operator-concave $f$ and a
positive unital map $T$ one has $T(fA)\le f(TA)$, with the subunital version
obtained by dilating $T$ to a unital map.  Mathlib~4.31 contains no such
theorem for positive maps.  Thus the three axioms are blocked on the
positive-map Jensen step, not on the integral representation of the integrand.

\subsection{The faithfulness caveat: completely positive versus positive}

The partial scaffold already present in the repository for the concave power
case\footnote{%
    \doclink{TNLean/Channel/Schwarz/OperatorJensenAux} (finite-POVM compression,
    Schur-complement inverse bound, and the resolvent inequality).}
proves the Jensen step through a Kraus / Stinespring dilation of $T$.  That
construction is available only for \emph{completely positive} maps: a positive
map that is not completely positive has no Kraus form and no Stinespring
dilation, so the compression argument does not apply to it.

Consequently a proof of the three axioms \emph{via that scaffold} would
establish the inequalities only under the added hypothesis that $T$ is
completely positive.  The axioms as stated assume only that $T$ is positive
(the hypothesis recorded in the formal statements is \leanid{IsPositiveMap}).
Discharging them through the dilation route while keeping the present signatures
would therefore prove a different, stronger-hypothesis statement than the one
asserted, in violation of the faithfulness rule.  Two faithful options remain:

\begin{enumerate}
  \item restate the three axioms with an explicit complete-positivity
    hypothesis, matching what the dilation argument proves and what Wolf's
    Corollary~5.2 in fact uses (the maps there are channels, hence completely
    positive); or
  \item prove the Choi--Davis--Jensen inequality by a positivity-only argument,
    keeping the bare \leanid{IsPositiveMap} hypothesis.  In the matrix-algebra
    setting the positive-map operator Jensen inequality $f(TA)\le T(fA)$ for an
    operator-convex $f$ and a \emph{unital positive} map $T$ is the classical
    Davis (1957) / Choi (1974) theorem, whose standard proof needs only
    positivity and unitality of $T$ (not complete positivity, and not the
    stronger $2$-positivity).  The subunital case is recovered by dilating $T$
    to a unital positive map.  This route avoids the Kraus/Stinespring
    construction entirely and so is faithful to the bare positivity hypothesis.
\end{enumerate}

\subsection{The concrete first step}\label{ssec:first-step}

Independently of the dilation question, one input is still missing from the
formal library on the convex side: operator convexity of $x\mapsto x^p$ for
$p\in[1,2]$.  Adding \leanid{CFC.convexOn_rpow} for $p\in[1,2]$ discharges the
corresponding upstream TODO\footnote{%
    \path{Mathlib/Analysis/SpecialFunctions/ContinuousFunctionalCalculus/Rpow/Order.lean},
    line~29.}
and supplies the scalar operator-convexity input for
\leanid{posMap_rpow_convex_jensen}.  It is obtained from the $[0,1]$ data
already in the library through the Bendat--Sherman / Hansen--Pedersen
characterization: for a continuous $f$ on $[0,\infty)$ with $f(0)\le 0$, $f$ is
operator convex if and only if $x\mapsto f(x)/x$ is operator monotone.  Applied
to $f(x)=x^p$ with $f(0)=0$, this gives: $x\mapsto x^p$ is operator convex on
$[1,2]$ if and only if $x\mapsto x^{p-1}$ is operator monotone on $[0,1]$, and
the latter is exactly \leanid{CFC.monotone_rpow} (operator monotonicity of
$x^{p-1}$ for $p-1\in[0,1]$).  The bridge therefore runs through operator
\emph{monotonicity}, not operator concavity: the naive identity
$x^p = x\cdot x^{p-1}$ does \emph{not} establish operator convexity, since
operator convexity is not preserved under multiplication by $x$.  This is the
smallest reusable increment toward the convex power case.

\subsection{Verdict}

The three positive-map Jensen statements are mathematically standard
\cite[Theorem~5.1]{Wolf2012Quantum}.  Their elimination is blocked on a
positive-map Choi--Davis--Jensen theorem, which is absent from Mathlib~4.31; the
integral representation of the integrand, previously named as the obstruction,
is present.  A proof through the existing dilation scaffold is faithful only
under an added complete-positivity hypothesis.

\section{The Lieb concavity axiom}

\subsection{The assertion}

For $s\in[0,1]$, an arbitrary matrix $K$, and positive-definite $A,B$, the map
\begin{align}
  (A,B)\ \longmapsto\ \tr\!\bigl(K^\dagger A^s K B^{1-s}\bigr)
  \tag{\path{lieb_concavity_axiom}}
\end{align}
is jointly concave on pairs of positive-definite matrices.  This is the Lieb
concavity theorem, the trace functional underlying the Wigner--Yanase--Dyson
conjecture.

\subsection{The point at issue}

The standard route to joint concavity is Ando's operator integral
representation
\begin{align}
  A^s B^{1-s} = \frac{\sin(\pi s)}{\pi}\int_0^\infty t^{s-1}\,
    A\,(A+tB)^{-1}\,B\ \mathrm dt,
\end{align}
combined with monotonicity of the resolvent $(A+tB)^{-1}$ in each argument.
Joint concavity of the integrand for each fixed $t$, followed by integration and
the linearity and positivity of the trace, yields the result.

Mathlib~4.31 contains neither this operator integral representation nor a
matrix-mean / resolvent-monotonicity API of the required form.  Supplying it is a
library-level contribution, not a short local bridge.

\subsection{Verdict}

The Lieb concavity statement is correct \cite{Wolf2012Quantum}.  Its elimination
requires the Ando integral representation of $A^s B^{1-s}$ and resolvent
monotonicity, both absent from Mathlib~4.31.  The axiom is therefore blocked at
the library level.  Three downstream corollaries already build on it (the
$K=\mathbf 1$ case and separate concavity in each argument); these inherit the
same blocked status until the integral representation is formalized.

\section*{References}

Beyond \cite{Wolf2012Quantum}, the underlying matrix-analysis facts are:
E.~H.\ Lieb, \emph{Convex trace functions and the Wigner--Yanase--Dyson
conjecture}, Adv.\ Math.\ \textbf{11} (1973) 267--288;
T.~Ando, \emph{Concavity of certain maps on positive definite matrices and
applications to Hadamard products}, Linear Algebra Appl.\ \textbf{26} (1979)
203--241;
F.~Hansen and G.~K.\ Pedersen, \emph{Jensen's operator inequality}, Bull.\
London Math.\ Soc.\ \textbf{35} (2003) 553--564;
R.~Bhatia, \emph{Matrix Analysis}, Springer GTM~169 (1997), Chapter~V.
The relevant formal-library lemmas are
\leanid{CFC.exists_measure_nnrpow_eq_integral_cfcₙ_rpowIntegrand₀₁},
\leanid{CFC.concaveOn_rpow}, and \leanid{CFC.concaveOn_log}; the missing
upstream item is \leanid{CFC.convexOn_rpow} for $p\in[1,2]$.

\bibliographystyle{alpha}
\bibliography{references}

\end{document}
